%% =====================================================================
%%  Supplementary material for the manuscript in main.tex.
%%
%%  Every number is a macro from macros_auto.tex and every table is generated
%%  by make_tables.py into supp_table_*_auto.tex, exactly as in the main text.
%%  Items are numbered S1, S2, ... and the main text cites them through the
%%  xr package, so build this file BEFORE main.tex:
%%      latexmk -pdf -outdir=build supplementary.tex
%%      latexmk -pdf -outdir=build main.tex
%% =====================================================================
\documentclass{iopjournal}

\usepackage{graphicx}
\usepackage{amsmath,amssymb}
\graphicspath{{../results/figures_paper/}{../results/figures/}{figures/}}
%% AUTO-GENERATED by paper/make_tables.py: DO NOT EDIT BY HAND.
%% Regenerate after every benchmark run:  python paper/make_tables.py


%% -------------------------------------------- inline macros

\newcommand{\NSubjects}{30}
\newcommand{\NPipelines}{15}
\newcommand{\NFoldScores}{6\,750}
\newcommand{\NTrials}{45}
\newcommand{\BestClassicalName}{classical/CSP+LDA}
\newcommand{\BestClassicalAcc}{0.611}
\newcommand{\BestQuantumName}{quantum/CNOT-kernel-SVM}
\newcommand{\BestQuantumAcc}{0.579}
\newcommand{\PrimaryDelta}{+0.032}
\newcommand{\PrimaryP}{\ensuremath{{}={}}0.002}
\newcommand{\PrimaryDz}{0.58}
\newcommand{\PrimaryBetter}{24/30}
\newcommand{\AblationDelta}{+0.0106}
\newcommand{\AblationP}{\ensuremath{{}={}}0.309}
\newcommand{\AblationBetter}{17/30}
\newcommand{\DimDelta}{+0.0163}
\newcommand{\DimP}{\ensuremath{{}={}}0.173}
\newcommand{\NSurviveHolm}{0}
\newcommand{\NFamilyTests}{14}
\newcommand{\FastestName}{classical/logvar+LDA}
\newcommand{\FastestSec}{0.046}
\newcommand{\SlowestName}{quantum/Bures-RBF-SVM}
\newcommand{\SlowestSec}{12.4}
\newcommand{\SpeedRatio}{272}

\newcommand{\ConcNSubjects}{10}

%% -------------------------------- cross-dataset macros

\newcommand{\BciNSubjects}{9}
\newcommand{\BciNTrials}{288}
\newcommand{\BciNFoldScores}{2\,025}
\newcommand{\BciBestClassicalAcc}{0.763}
\newcommand{\BciBestQuantumAcc}{0.685}
\newcommand{\BciWorstQuantumAcc}{0.580}
\newcommand{\GapPhysio}{0.032}
\newcommand{\GapBci}{0.079}
\newcommand{\GapWorstPhysio}{0.088}
\newcommand{\GapWorstBci}{0.183}
\newcommand{\SpearmanRho}{0.882}
\newcommand{\BciPrimaryDelta}{+0.066}
\newcommand{\BciPrimaryP}{\ensuremath{{}={}}0.004}
\newcommand{\BciPrimaryDz}{1.14}
\newcommand{\BciPrimaryBetter}{9/9}
\newcommand{\BciDensDelta}{+0.171}
\newcommand{\BciDensDz}{1.45}
\newcommand{\BciDensBetter}{9/9}
\newcommand{\FisherPrimaryP}{\ensuremath{{}<{}}0.001}
\newcommand{\FisherDensP}{\ensuremath{{}<{}}0.001}
\newcommand{\FisherAblationP}{\ensuremath{{}={}}0.237}
\newcommand{\BciWilcoxonFloor}{0.0039}
\newcommand{\BciHolmFloor}{0.0547}


%% ------------------------------ reference-frame macros

\newcommand{\NPipelinesExt}{23}
\newcommand{\NPipelinesFb}{15}
\newcommand{\PhysFrameMin}{+0.050}
\newcommand{\PhysFrameMax}{+0.106}
\newcommand{\PhysFrameBestKernel}{Fidelity}
\newcommand{\PhysFrameMaxP}{0.049}
\newcommand{\PhysFrameAllBetter}{no}
\newcommand{\PhysFrameN}{30}
\newcommand{\TwinName}{control/riemann-kernel-SVM}
\newcommand{\TwinAcc}{0.621}
\newcommand{\TwinDeltaMin}{-0.0151}
\newcommand{\TwinDeltaMax}{+0.0081}
\newcommand{\TwinMinP}{\ensuremath{{}={}}0.183}
\newcommand{\HeadBestClassical}{classical/CSP+LDA}
\newcommand{\HeadSensorDelta}{+0.0523}
\newcommand{\HeadSensorP}{\ensuremath{{}={}}0.016}
\newcommand{\HeadSensorBetter}{20/30}
\newcommand{\HeadRefDelta}{-0.0183}
\newcommand{\HeadRefLead}{0.0183}
\newcommand{\HeadRefP}{\ensuremath{{}={}}0.059}
\newcommand{\HeadRefBetter}{10/30}
\newcommand{\HeadRefKernel}{quantum/Fidelity-ref-SVM}
\newcommand{\BciFrameMin}{+0.164}
\newcommand{\BciFrameMax}{+0.187}
\newcommand{\BciFrameBestKernel}{Fidelity}
\newcommand{\BciFrameMaxP}{0.004}
\newcommand{\BciFrameAllBetter}{yes}
\newcommand{\BciFrameN}{9}
\newcommand{\FourQFrameMin}{+0.075}
\newcommand{\FourQFrameMax}{+0.106}
\newcommand{\FourQFrameBestKernel}{HS overlap}
\newcommand{\FourQFrameMaxP}{0.004}
\newcommand{\FourQFrameAllBetter}{no}
\newcommand{\FourQFrameN}{30}
\newcommand{\FourQTwinName}{control/riemann-kernel-SVM}
\newcommand{\FourQTwinAcc}{0.625}
\newcommand{\FourQBestKernel}{quantum/Bures-RBF-ref-SVM}
\newcommand{\FourQBestAcc}{0.626}
\newcommand{\FourQTwinDeltaMin}{-0.0114}
\newcommand{\FourQTwinDeltaMax}{+0.0005}
\newcommand{\FourQTwinMinP}{\ensuremath{{}={}}0.253}
\newcommand{\FourQNSubjects}{30}
\newcommand{\FourQTwinTopKernel}{quantum/Bures-RBF-ref-SVM}
\newcommand{\FourQTwinTopDelta}{+0.0005}
\newcommand{\FourQTwinTopP}{\ensuremath{{}={}}0.650}
\newcommand{\FourQTwinTopBetter}{14/30}
\newcommand{\FourQTwinRestDeltaMin}{-0.0114}
\newcommand{\FourQTwinRestDeltaMax}{-0.0015}
\newcommand{\FourQTwinRestMinP}{\ensuremath{{}={}}0.253}
\newcommand{\FourQBestClassical}{classical/TS+LR}
\newcommand{\FourQBestClassicalAcc}{0.606}
\newcommand{\FourQHeadRefDelta}{+0.0200}
\newcommand{\FourQHeadRefP}{\ensuremath{{}={}}0.016}
\newcommand{\FourQHeadRefBetter}{21/30}
\newcommand{\FourQTwinHeadDelta}{+0.0195}
\newcommand{\FourQTwinHeadP}{\ensuremath{{}={}}0.017}
\newcommand{\FourQTwinHeadBetter}{17/30}
\newcommand{\ChoFrameMin}{+0.012}
\newcommand{\ChoFrameMax}{+0.063}
\newcommand{\ChoFrameBestKernel}{HS overlap}
\newcommand{\ChoFrameMaxP}{0.009}
\newcommand{\ChoFrameAllBetter}{no}
\newcommand{\ChoFrameN}{52}
\newcommand{\ChoTwinName}{control/riemann-kernel-SVM}
\newcommand{\ChoTwinAcc}{0.638}
\newcommand{\ChoBestKernel}{quantum/Fidelity-ref-SVM}
\newcommand{\ChoBestAcc}{0.648}
\newcommand{\ChoTwinDeltaMin}{-0.0008}
\newcommand{\ChoTwinDeltaMax}{+0.0100}
\newcommand{\ChoTwinMinP}{\ensuremath{{}={}}0.001}
\newcommand{\ChoNSubjects}{52}
\newcommand{\ChoTwinTopKernel}{quantum/Fidelity-ref-SVM}
\newcommand{\ChoTwinTopDelta}{+0.0100}
\newcommand{\ChoTwinTopP}{\ensuremath{{}={}}0.001}
\newcommand{\ChoTwinTopBetter}{39/52}
\newcommand{\ChoTwinRestDeltaMin}{-0.0008}
\newcommand{\ChoTwinRestDeltaMax}{+0.0041}
\newcommand{\ChoTwinRestMinP}{\ensuremath{{}={}}0.128}
\newcommand{\ChoBestClassical}{classical/TS+LR}
\newcommand{\ChoBestClassicalAcc}{0.644}
\newcommand{\ChoHeadRefDelta}{+0.0039}
\newcommand{\ChoHeadRefP}{\ensuremath{{}={}}0.174}
\newcommand{\ChoHeadRefBetter}{31/52}
\newcommand{\ChoTwinHeadDelta}{-0.0061}
\newcommand{\ChoTwinHeadP}{\ensuremath{{}={}}0.010}
\newcommand{\ChoTwinHeadBetter}{18/52}
\newcommand{\ChoNTrialsMin}{200}
\newcommand{\ChoNTrialsMax}{240}
\newcommand{\ChoNTrialsMaxCount}{3}
\newcommand{\ChoHeadSensorDelta}{+0.0152}
\newcommand{\ChoHeadSensorP}{\ensuremath{{}<{}}0.001}
\newcommand{\ChoSensorBestQuantum}{quantum/QRE-RBF-SVM}
\newcommand{\SeedRbfWorst}{-0.0284}
\newcommand{\SeedRbfMinP}{\ensuremath{{}<{}}0.001}
\newcommand{\SeedOverlapDeltaMin}{-0.0079}
\newcommand{\SeedOverlapDeltaMax}{+0.0081}
\newcommand{\SeedOverlapMinP}{\ensuremath{{}={}}0.194}
\newcommand{\SeedSigCount}{5}
\newcommand{\SeedSigN}{15}
\newcommand{\SeedSigFavourQuantum}{0}
\newcommand{\SeedSigFavourTwin}{5}
\newcommand{\SeedN}{3}
\newcommand{\SeedFrameMin}{+0.032}
\newcommand{\SeedFrameMax}{+0.107}
\newcommand{\SeedTwinDeltaMin}{-0.0284}
\newcommand{\SeedTwinDeltaMax}{+0.0081}
\newcommand{\SeedTwinMinP}{\ensuremath{{}<{}}0.001}
\newcommand{\SeedEquivPass}{4}
\newcommand{\SeedEquivN}{15}
\newcommand{\SeedWorstBound}{0.042}
\newcommand{\SeedHeadRefMin}{+0.0183}
\newcommand{\SeedHeadRefMax}{+0.0215}
\newcommand{\SeedHeadReversals}{3}
\newcommand{\PhysCostTwinSec}{3.2}
\newcommand{\PhysCostSlowName}{quantum/Bures-RBF-ref-SVM}
\newcommand{\PhysCostSlowSec}{12.1}
\newcommand{\PhysCostRatio}{3.8}
\newcommand{\BciCostTwinSec}{15.7}
\newcommand{\BciCostSlowName}{quantum/Bures-RBF-ref-SVM}
\newcommand{\BciCostSlowSec}{117.7}
\newcommand{\BciCostRatio}{7.5}
\newcommand{\FbFrameMin}{+0.124}
\newcommand{\FbFrameMax}{+0.152}
\newcommand{\FbFrameBestKernel}{Fidelity}
\newcommand{\FbFrameMaxP}{\ensuremath{<}0.001}
\newcommand{\FbFrameAllBetter}{no}
\newcommand{\FbFrameN}{30}
\newcommand{\FbcspDelta}{+0.0212}
\newcommand{\FbcspP}{\ensuremath{{}={}}0.125}
\newcommand{\FbcspBestAcc}{0.648}
\newcommand{\FbcspAcc}{0.627}
\newcommand{\TransferN}{30}
\newcommand{\TransferSpread}{0.0141}
\newcommand{\TransferBest}{quantum/HS-overlap}
\newcommand{\TransferBestAcc}{0.647}
\newcommand{\TransferFloor}{1.9\times10^{-9}}
\newcommand{\CsN}{9}
\newcommand{\CsFloor}{3.9\times10^{-3}}
\newcommand{\CsTrialsPerSession}{144}
\newcommand{\CsFrameQMin}{+0.169}
\newcommand{\CsFrameQMax}{+0.217}
\newcommand{\CsFrameQMaxP}{0.008}
\newcommand{\CsFrameQAllBetter}{no}
\newcommand{\CsSensorQMin}{0.542}
\newcommand{\CsSensorQMax}{0.578}
\newcommand{\CsRefQMin}{0.746}
\newcommand{\CsRefQMax}{0.760}
\newcommand{\CsBestKernel}{quantum/Bures-RBF}
\newcommand{\CsBestAcc}{0.760}
\newcommand{\CsFrameClMin}{+0.022}
\newcommand{\CsFrameClMax}{+0.031}
\newcommand{\CsSensorClMin}{0.701}
\newcommand{\CsSensorClMax}{0.717}
\newcommand{\CsBestClassical}{classical/TS+LR}
\newcommand{\CsBestClassicalAcc}{0.738}
\newcommand{\CsTwinName}{control/riemann-kernel-SVM}
\newcommand{\CsTwinAcc}{0.755}
\newcommand{\CsTwinDeltaMin}{-0.0085}
\newcommand{\CsTwinDeltaMax}{+0.0054}
\newcommand{\CsTwinMinP}{\ensuremath{{}={}}0.500}
\newcommand{\CsTwinAbsMax}{0.0085}
\newcommand{\CsSpread}{0.0139}
\newcommand{\SweepQMin}{3}
\newcommand{\SweepQMax}{6}
\newcommand{\SweepN}{14}
\newcommand{\SweepRefRatioMin}{1.6}
\newcommand{\SweepRefRatioMax}{2.4}
\newcommand{\SweepSensorRatioMin}{4.1}
\newcommand{\SweepSensorRatioMax}{7.7}
\newcommand{\SweepGainLoMin}{4.7}
\newcommand{\SweepGainLoMax}{9.5}
\newcommand{\SweepGainHiMin}{1.9}
\newcommand{\SweepGainHiMax}{2.6}
\newcommand{\SweepRefAboveBase}{every}
\newcommand{\SweepRefMonotone}{not every}
\newcommand{\SweepRefDipMax}{8}
\newcommand{\GramGainMin}{4.7}
\newcommand{\GramGainMax}{9.5}
\newcommand{\GramNSubjects}{14}
\newcommand{\GramSensorMean}{0.981}
\newcommand{\GramSensorStd}{0.022}
\newcommand{\ShotCrossover}{10^{4}}
\newcommand{\FrameEffectMin}{+0.012}
\newcommand{\FrameEffectMax}{+0.217}


%% ------------------------------ equivalence macros

\newcommand{\EquivMargin}{0.02}
\newcommand{\EquivN}{35}
\newcommand{\EquivPass}{20}
\newcommand{\EquivWorstBound}{0.032}
\newcommand{\EquivFailQuantum}{10}
\newcommand{\EquivFailClassical}{5}
\newcommand{\EquivNZeroExcluded}{3}
\newcommand{\EquivNSettings}{seven}
\newcommand{\EquivNDatasets}{three}
\newcommand{\EquivNRegisters}{three}
\newcommand{\EquivBciPass}{5}
\newcommand{\EquivBciN}{5}
\newcommand{\EquivBciWorst}{0.0109}
\newcommand{\EquivCsPass}{2}
\newcommand{\EquivCsN}{5}
\newcommand{\EquivCsWorst}{0.0281}
\newcommand{\EquivTrPass}{3}
\newcommand{\EquivTrN}{5}
\newcommand{\EquivTrWorst}{0.0235}
\newcommand{\EquivFourQPass}{4}
\newcommand{\EquivFourQN}{5}
\newcommand{\EquivFourQWorst}{0.0263}
\newcommand{\EquivChoPass}{5}
\newcommand{\EquivChoN}{5}
\newcommand{\EquivChoWorst}{0.0153}
\newcommand{\EquivZeroExclName}{Fidelity (PhysioNet, 5 qubits)}
\newcommand{\EquivZeroExclDelta}{+0.0180}
\newcommand{\EquivZeroExclList}{Fidelity on PhysioNet, 5 qubits at $+0.0180$, HS overlap on Transfer (LOSO) at $+0.0119$ and Fidelity on Cho2017 at $+0.0100$}
\newcommand{\EquivZeroExclAllQuantum}{yes}


\renewcommand{\thetable}{S\arabic{table}}
\renewcommand{\thefigure}{S\arabic{figure}}
\renewcommand{\thesection}{S\arabic{section}}

\begin{document}

\articletype{Supplementary material}

\title{Supplementary material for: Quantum Kernels for EEG: An Invariance
Mismatch Masquerading as Advantage}

\author{Md.~Imtiaj~Alam~Sajin$^{1,*}$\orcid{0009-0009-2423-1835},
Md~Wahiduzzaman~Suva$^{1}$,
Esm~E~Moula~Chowdhury~Abha$^{1}$,
Samiul~Islam~Saif$^{1}$\orcid{0009-0009-7116-7151} and
Muhammad~Hasibur~Rashid~Chayon$^{1}$\orcid{0000-0003-3741-6651}}

\affil{$^1$Department of Computer Science, American International
University-Bangladesh (AIUB), Dhaka, Bangladesh}

\affil{$^*$Author to whom any correspondence should be addressed.}

\email{26-94090-2@student.aiub.edu}

\keywords{supplementary tables and figures}

\begin{abstract}
Hyperparameter grids for every pipeline, the per-subject differences behind
the main benchmark, the kernel-concentration and register-sweep tables, the
cross-session and register-sweep figure, and the cross-subject and finite-shot
figure and table. Each item is referred to from the main text.
\end{abstract}

\section{Hyperparameter grids}

%% AUTO-GENERATED by paper/make_tables.py: DO NOT EDIT BY HAND.
%% Regenerate after every benchmark run:  python paper/make_tables.py


%% --------------------------------------------------- Table: grids
\begin{table}[htbp]
\caption{\label{tab:grids}Hyperparameter grids and fixed settings. Every grid is
searched by the same inner cross-validation (4-fold
within subject; 3-fold subject-grouped for
cross-subject transfer; 4-fold stratified
within the training session for cross-session transfer). The grids and
settings are read from the code that produced the results. $\gamma_0$ is the
median-heuristic bandwidth, computed on the training split only.}

\begin{tabular}{@{}p{0.40\textwidth}p{0.55\textwidth}@{}}
\hline
Pipelines & Searched grid \\
\hline
FB-Fidelity-SVM, FB-Fidelity-ref-SVM, FB-HS-overlap-SVM, FB-HS-overlap-ref-SVM, FB-TS+LR, FB-logeuclid-kernel-SVM, FB-riemann-kernel-SVM, Fidelity-SVM, Fidelity-ref-SVM, HS-overlap-SVM, HS-overlap-ref-SVM, PCA-matched-linear, TS+LR, logeuclid-TS+LR, logeuclid-kernel-SVM, riemann-kernel-SVM & $C$ $\{$0.1, 1, 10$\}$ \\
FB-TS+RBF-SVM, PCA-matched-RBF, TS+RBF-SVM & $C$ $\{$0.1, 1, 10$\}$; RBF $\gamma$ $\{$\texttt{scale}, 0.01, 0.1$\}$ \\
CNOT-kernel-SVM, IQP-kernel-SVM, IQP-no-entangle & $C$ $\{$0.1, 1, 10$\}$; angle scale $\{$0.5, 1, 2$\}$ \\
Bures-RBF-SVM, Bures-RBF-ref-SVM, FB-Bures-RBF-SVM, FB-Bures-RBF-ref-SVM, FB-HS-RBF-SVM, FB-HS-RBF-ref-SVM, FB-QRE-RBF-SVM, FB-QRE-RBF-ref-SVM, HS-RBF-SVM, HS-RBF-ref-SVM, QRE-RBF-SVM, QRE-RBF-ref-SVM & $C$ $\{$0.1, 1, 10$\}$; bandwidth multiplier $\times\gamma_0$ $\{$0.25, 1, 4$\}$ \\
CSP+LDA & CSP filters $\{$2, 4, 6$\}$ \\
FBCSP+LDA & CSP filters per band $\{$2, 4, 6$\}$ \\
MDM, logvar+LDA & none (no tuned hyperparameter) \\
\hline
Cross-subject and cross-session quantum kernels & $C$ $\{$0.1, 1, 10$\}$; bandwidth multiplier $\times\gamma_0$ $\{$0.25, 1, 4$\}$ for distance kernels \\
\hline
\multicolumn{2}{@{}p{0.95\textwidth}@{}}{\emph{Fixed in every pipeline:} OAS shrinkage covariance; LDA with
Ledoit--Wolf shrinkage; tangent space at the Riemannian mean with feature
standardisation; circuit kernels and dimension-matched controls on a PCA
reduction to 4 components scaled to $[0,\pi]$, 2 circuit layers; SVMs on
precomputed kernels projected to the positive semi-definite cone; logistic
regression with up to 2000 iterations.} \\
\hline
\end{tabular}
\end{table}


\section{Per-subject differences in the sensor-frame benchmark}

\begin{figure}[htbp]
\centering
\includegraphics[width=\textwidth]{fig3_paired_differences.pdf}
\caption{\label{fig:paired}Paired per-subject accuracy differences against the
Riemannian tangent-space reference. Dots are individual subjects; diamonds are
means. Points left of zero indicate worse performance than the reference.}
\end{figure}

\section{Kernel concentration versus register size}

%% AUTO-GENERATED by paper/make_tables.py: DO NOT EDIT BY HAND.
%% Regenerate after every benchmark run:  python paper/make_tables.py


%% ---------------------------------------------------------------- Table 4
\begin{table}[htbp]
\caption{\label{tab:concentration}Kernel concentration as a function of register
size on real EEG. Register size is swept by varying the channel count over
powers of two (4/8/16/32/64 channels $=$ 2--6 qubits), with no change of method.
The statistic is the variance of the off-diagonal Gram entries; a factor below
one means concentration worsens with scale. The entangled circuit kernel concentrates $2.7\times$ faster in log-slope than the same circuit with entanglers removed.}
\begin{tabular}{@{}lcccc@{}}
\hline
Kernel & Log-slope / qubit & Factor / qubit & Var (2 qubits) & Var (6 qubits) \\
\hline
IQP-entangled & $-0.859$ & 0.424 & 0.07476 & 0.00248 \\
IQP-product & $-0.319$ & 0.727 & 0.09396 & 0.02652 \\
Bures-RBF & $-0.195$ & 0.823 & 0.07758 & 0.03643 \\
HS-RBF & $-0.113$ & 0.894 & 0.07682 & 0.05009 \\
Fidelity & $+0.493$ & 1.638 & 0.00030 & 0.00213 \\
HS-overlap & $+0.536$ & 1.709 & 0.00043 & 0.00369 \\
\hline
\end{tabular}
\end{table}


%% AUTO-GENERATED by paper/make_tables.py: DO NOT EDIT BY HAND.
%% Regenerate after every benchmark run:  python paper/make_tables.py


%% ------------------------------------------------------ Table: sweep
\begin{table}[htbp]
\caption{\label{tab:sweep}Off-diagonal Gram variance versus register size in
both frames, $n=14$ subjects, PhysioNet.
Larger is better: a kernel whose Gram variance collapses cannot separate
trials. ``Gain'' is reference over sensor at the same register; ``vs.\ 3\,q''
is the reference-frame variance relative to its own three-qubit value. In the
reference frame no kernel's variance falls below its three-qubit value at any
larger register, so the concentration that the sensor frame exhibits is not a
qubit-count effect once the frame is corrected.}
\begin{tabular}{@{}llcccc@{}}
\hline
Register & Kernel & Sensor ($\times10^{-3}$) & Reference ($\times10^{-3}$) &
Gain & vs.\ 3\,q \\
\hline
3\,q (8 ch) & Bures & 0.35 & 1.78 & 5.1$\times$ & 1.00 \\
 & Fidelity & 0.33 & 1.55 & 4.7$\times$ & 1.00 \\
 & HS overlap & 0.49 & 4.71 & 9.5$\times$ & 1.00 \\
 & QRE & 1.58 & 8.27 & 5.2$\times$ & 1.00 \\
\hline
4\,q (16 ch) & Bures & 0.58 & 3.20 & 5.6$\times$ & 1.80 \\
 & Fidelity & 0.54 & 2.46 & 4.6$\times$ & 1.59 \\
 & HS overlap & 0.73 & 6.56 & 9.0$\times$ & 1.39 \\
 & QRE & 2.84 & 15.72 & 5.5$\times$ & 1.90 \\
\hline
5\,q (32 ch) & Bures & 0.69 & 3.72 & 5.4$\times$ & 2.09 \\
 & Fidelity & 0.64 & 2.73 & 4.3$\times$ & 1.76 \\
 & HS overlap & 0.94 & 7.77 & 8.3$\times$ & 1.65 \\
 & QRE & 3.43 & 20.23 & 5.9$\times$ & 2.45 \\
\hline
6\,q (64 ch) & Bures & 1.54 & 3.44 & 2.2$\times$ & 1.93 \\
 & Fidelity & 1.36 & 2.55 & 1.9$\times$ & 1.65 \\
 & HS overlap & 3.83 & 9.01 & 2.4$\times$ & 1.91 \\
 & QRE & 7.78 & 20.11 & 2.6$\times$ & 2.43 \\
\hline
\end{tabular}
\end{table}


\section{Cross-session transfer and the register sweep}

\begin{figure}[htbp]
\centering
\includegraphics[width=\textwidth]{fig10_crosssession_sweep.pdf}
\caption{\label{fig:crosssession}(a) Cross-session transfer within subject on
IV-2a, both directions averaged. In the sensor frame the quantum kernels sit
near chance while the classical baselines do not; whitening each session by
its own label-free mean closes the gap and leaves every SPD-kernel geometry
tied. (b) Off-diagonal Gram variance versus register size in both frames on
PhysioNet ($n=\SweepN$). Solid, reference frame; dashed, sensor frame. In the
reference frame no kernel falls below its three-qubit value at any larger
register.}
\end{figure}

\section{Cross-subject transfer and finite-shot estimation}

\begin{figure}[htbp]
\centering
\includegraphics[width=\textwidth]{fig9_transfer_shots.pdf}
\caption{\label{fig:transfer}(a) Cross-subject transfer, leave-one-subject-out.
Recentring helps every method and helps the quantum kernels (Q) most, but in
the reference frame all nine geometries are indistinguishable. (b) Accuracy
against shot budget. The reference frame needs more shots to approach its own
ceiling because that ceiling is higher, yet from $\ShotCrossover$ shots per
Gram entry it already exceeds what the sensor frame attains with unlimited
shots.}
\end{figure}

%% AUTO-GENERATED by paper/make_tables.py: DO NOT EDIT BY HAND.
%% Regenerate after every benchmark run:  python paper/make_tables.py


%% ------------------------------------------------------- Table: shot budget
\begin{table}[htbp]
\caption{\label{tab:shots}Accuracy under finite-shot estimation. Unlike every
other table, the reference state here is estimated from all of a subject's
trials, not the training split alone; it uses no labels, but these accuracies
are comparable only with one another, not with the nested cross-validation
results. $\mathrm{tr}(\rho\sigma)$ is the SWAP-test observable, so $S$ shots give an
unbiased estimate with variance $(1-k^2)/S$; we sample each unordered pair
binomially and project the Gram matrix back to the positive semi-definite cone.
The reference frame needs \emph{more} shots to approach its own ceiling,
because that ceiling is higher, but it dominates the sensor frame in absolute
terms from $10^4$ shots upwards, above which it exceeds what the sensor
frame achieves with unlimited shots.}
\begin{tabular}{@{}llcccccc@{}}
\hline
Kernel & Frame & $10^{2}$ & $10^{3}$ & $10^{4}$ & $10^{5}$ & $10^{6}$ & $\infty$ \\
\hline
HS-RBF & Sensor & 0.492 & 0.505 & 0.516 & 0.539 & 0.546 & \textbf{0.549} \\
HS-RBF & Reference & 0.508 & 0.549 & 0.580 & 0.614 & 0.620 & \textbf{0.627} \\
HS-overlap & Sensor & 0.498 & 0.509 & 0.519 & 0.530 & 0.528 & \textbf{0.528} \\
HS-overlap & Reference & 0.521 & 0.542 & 0.585 & 0.615 & 0.628 & \textbf{0.631} \\
\hline
\end{tabular}
\end{table}


\section{Few-trial calibration}

\begin{figure}[htbp]
\centering
\includegraphics[width=\textwidth]{fig11_calibration.pdf}
\caption{\label{fig:calibration}Accuracy against the number of calibration
trials, and the quantum minus twin difference that goes with it. The
reference-frame kernels hold a small lead over the metric-matched twin while
calibration trials are scarce, and it closes as trials are added. The main
text gives the sizes, their corrections,
and the limits that bound the result: its mean lies inside the equivalence
margin, it clears the best classical pipeline only at the smallest training
sets, and the pre-specified trend test does not survive correction.}
\end{figure}

\end{document}
